% Roteiro de fala — Projeto ECOX14 (Tópicos em Programação, UNIFEI)
% Compilar com: xelatex roteiro_apresentacao.tex
\documentclass[12pt,a4paper]{article}

\usepackage[brazil]{babel}
\usepackage{fontspec}
\setmainfont{Carlito}
\setmonofont{DejaVu Sans Mono}[Scale=0.85]
\usepackage[a4paper,margin=2.2cm,top=2cm,bottom=2.2cm,headheight=15pt]{geometry}
\usepackage{xcolor}
\usepackage[most]{tcolorbox}
\usepackage{booktabs}
\usepackage{tabularx}
\usepackage{enumitem}
\usepackage{titlesec}
\usepackage{fancyhdr}
\usepackage{microtype}
\usepackage[hidelinks]{hyperref}

% ---------- Cores ----------
\definecolor{azul}{HTML}{2A78D6}
\definecolor{azulclaro}{HTML}{EAF2FC}
\definecolor{laranja}{HTML}{EB6834}
\definecolor{laranjaclaro}{HTML}{FDF0EA}
\definecolor{cinza}{HTML}{52514E}
\definecolor{cinzaclaro}{HTML}{F3F2EE}

% ---------- Títulos ----------
\titleformat{\section}{\Large\bfseries}{\thesection}{0.6em}{}
\titleformat{\subsection}{\large\bfseries}{}{0em}{}
\titlespacing*{\section}{0pt}{1.4em}{0.6em}
\titlespacing*{\subsection}{0pt}{1.2em}{0.4em}
\setlength{\parindent}{0pt}
\setlength{\parskip}{0.5em}
\setlist{itemsep=0.25em, topsep=0.3em}

% ---------- Cabeçalho ----------
\pagestyle{fancy}
\fancyhf{}
\fancyhead[L]{\small\color{cinza} Roteiro de apresentação — ECOX14}
\fancyhead[R]{\small\color{cinza} Thiago}
\fancyfoot[C]{\small\color{cinza}\thepage}
\renewcommand{\headrulewidth}{0.4pt}

% ---------- Caixas ----------
% Texto para falar em voz alta
\newtcolorbox{fala}{enhanced, breakable, colback=azulclaro, colframe=azul,
  boxrule=0pt, leftrule=3pt, arc=2pt, left=10pt, right=8pt, top=6pt, bottom=6pt,
  fontupper=\large, before upper={\setlength{\parskip}{0.5em}}}
% O que fazer na tela
\newcommand{\natela}[1]{\par{\small\color{cinza}\textbf{Na tela:} #1}\par}
% Dica ou cuidado
\newtcolorbox{cuidado}{enhanced, breakable, colback=laranjaclaro, colframe=laranja,
  boxrule=0pt, leftrule=3pt, arc=2pt, left=10pt, right=8pt, top=5pt, bottom=5pt}
% Cabeçalho de cada tela: tempo + número + título
\newcommand{\tela}[3]{\subsection*{\colorbox{cinzaclaro}{\strut\texttt{#1}}\quad Tela #2 — #3}}

\begin{document}

% ======================================================================
\begin{center}
{\LARGE\bfseries Roteiro de apresentação}\\[0.4em]
{\large Projeto de Engenharia de Dados — Tópicos em Programação (ECOX14)}\\[0.2em]
{\color{cinza} UNIFEI · tempo total: 10 minutos · apoio: \texttt{apresentacao.html}}
\end{center}

\section*{Como usar este roteiro}
\begin{itemize}
  \item Abra \texttt{apresentacao.html} no navegador (funciona sem internet). As setas \texttt{←} e \texttt{→} trocam de tela, a tecla \texttt{T} alterna entre tema claro e escuro, e os pontinhos no rodapé levam direto a qualquer tela.
  \item O texto nas caixas azuis é para falar em voz alta. Não precisa decorar: leia uma vez em voz alta antes e fale com as suas palavras.
  \item O tempo no início de cada tela é o \emph{relógio acumulado}. Se ao terminar a tela 6 o relógio já passou de 5:40, encurte a tela 8.
\end{itemize}

% ======================================================================
\section{O que é importante falar}

Se a professora lembrar de só cinco coisas, que sejam estas. Todo o resto do roteiro serve para levar até elas.

\begin{enumerate}[leftmargin=*, label=\textbf{\arabic*.}]
  \item \textbf{A pergunta norteadora decide o que entra no projeto.} Ela não é enfeite: é o critério para criar (ou não) uma coluna. Foi ela que criou \texttt{calorias\_por\_hora} e que descartou recriar o IMC.
  \item \textbf{O pipeline é reprodutível e rastreável.} A bronze guarda o CSV como veio, a prata guarda o dado limpo, e cada execução grava o que foi decidido e com que números.
  \item \textbf{Limpar é medir, não supor, e dar o tipo certo a cada coluna.} A base veio sem ausentes e sem repetidas, e isso foi conferido. Sexo e tipo de treino viraram categoria; experiência virou categoria com ordem, porque 3 não é o triplo de 1.
  \item \textbf{O atributo derivado mudou a resposta.} Por sessão, o avançado queima 74\% a mais que o iniciante. Por hora, todos ficam em torno de 720 kcal/h: o avançado treina mais \emph{tempo}, não mais \emph{forte}.
  \item \textbf{Desconfiar de resultado bom demais.} Um R² de 98\% com quatro colunas levou a descobrir a fórmula que gerou a base. A resposta é honesta: vale para esta base, que é sintética, não para pessoas reais.
\end{enumerate}

\textbf{Onde não gastar tempo:} ler código linha por linha ou descrever as 15 colunas da base uma por uma. Se a professora quiser, ela pergunta.

% ======================================================================
\section{Roteiro tela a tela}

\begin{center}
\small
\begin{tabularx}{\linewidth}{@{}l l X r@{}}
\toprule
\textbf{Tela} & \textbf{Início} & \textbf{Assunto} & \textbf{Duração}\\
\midrule
1 & 0:00 & Abertura & 40 s\\
2 & 0:40 & Pergunta norteadora & 50 s\\
3 & 1:30 & A base de dados & 60 s\\
4 & 2:30 & Pipeline bronze → prata → ouro & 60 s\\
5 & 3:30 & Decisões de limpeza & 60 s\\
6 & 4:30 & Atributo derivado \texttt{calorias\_por\_hora} & 70 s\\
7 & 5:40 & O que explica as calorias & 60 s\\
8 & 6:40 & Explorador (pode encurtar) & 40 s\\
9 & 7:20 & A fórmula da base & 80 s\\
10 & 8:40 & Resposta, limitações e próximos passos & 80 s\\
\bottomrule
\end{tabularx}
\end{center}

% ----------------------------------------------------------------------
\tela{0:00}{1}{Abertura}
\begin{fala}
Bom dia. Meu projeto parte de uma pergunta simples: o que explica quantas calorias uma pessoa queima num treino? Para responder, montei um pipeline de dados em camadas sobre uma base pública de treinos em academia, com 973 praticantes e quatro tipos de treino: cardio, HIIT, força e yoga.
\end{fala}
\natela{só apresentar os quatro números e passar.}

% ----------------------------------------------------------------------
\tela{0:40}{2}{Pergunta norteadora}
\begin{fala}
A pergunta norteadora é: \emph{quais características do praticante mais explicam as calorias queimadas por sessão?}

Ela não serve só de título. Ela é o critério de decisão do projeto: se eu não consigo dizer em uma frase por que uma coluna ajuda a responder essa pergunta, a coluna não deveria existir. À direita estão as candidatas, ou seja, as colunas que a própria base traz: idade, sexo, peso, batimentos, duração, tipo de treino, experiência e outras.
\end{fala}

% ----------------------------------------------------------------------
\tela{1:30}{3}{A base de dados}
\begin{fala}
A base é a Gym Members Exercise, do Kaggle: um CSV baixado com token, com 973 linhas e 15 colunas. Cada linha é um praticante numa sessão de treino.

Ela é bem equilibrada: 511 homens e 462 mulheres, os quatro tipos de treino com 220 a 260 pessoas cada, e três níveis de experiência. À direita dá para ver a distribuição de cada coluna: as calorias por sessão vão de 303 a 1.783, com média de 905.

E guardem esta [\emph{trocar para Idade}]: a idade é quase plana de 18 a 59 anos. Vou voltar nisso no final.
\end{fala}
\natela{clique em ``Sexo'', ``Tipo de treino'' e ``Experiência'' enquanto fala da composição. No histograma, comece em ``Calorias por sessão'' e termine em ``Idade''.}

% ----------------------------------------------------------------------
\tela{2:30}{4}{Pipeline bronze → prata → ouro}
\begin{fala}
A arquitetura tem três camadas. Na bronze, o \texttt{ingerir.py} baixa o CSV do Kaggle e guarda exatamente como veio, com a data da extração no nome e um arquivo de proveniência dizendo de onde veio e quando. O \texttt{explorar.py} gera um relatório de perfilamento antes de qualquer decisão.

Na prata, o \texttt{transformar\_gym\_members.py} limpa, dá tipo às colunas, cria o atributo derivado e salva em Parquet: as mesmas 973 linhas, agora com 18 colunas. As funções genéricas de limpeza ficam num módulo separado, o \texttt{limpeza.py}, para não ficarem copiadas. E cada execução grava as decisões e os números reais daquela rodada. A ouro é a próxima etapa.
\end{fala}
\natela{clique em \texttt{ingerir.py} e depois em \texttt{transformar\_gym\_members.py}.}

% ----------------------------------------------------------------------
\tela{3:30}{5}{Decisões de limpeza}
\begin{fala}
A base veio limpa: zero valores ausentes e zero linhas repetidas. Mas ``veio limpa'' só vale se foi medido, e essas contagens ficam registradas. Se uma extração futura vier suja, o mesmo script avisa.

As decisões de verdade foram de tipo. Sexo e tipo de treino viraram categoria sem ordem: como texto livre, nada impede um ``Yogaa'' de entrar. O nível de experiência vem como 1, 2 e 3, mas 3 não é o triplo de 1: virou categoria com ordem, iniciante, intermediário e avançado. E criei uma chave de comparação do tipo de treino, sem acento e em minúscula, ao lado do rótulo original: uma serve para agrupar, a outra para exibir.
\end{fala}
\natela{abra os cartões na ordem enquanto fala. Se o tempo apertar, abra só o primeiro e o terceiro.}

% ----------------------------------------------------------------------
\tela{4:30}{6}{Atributo derivado \texttt{calorias\_por\_hora}}
\begin{fala}
O atributo derivado principal é calorias por hora: calorias dividido pela duração da sessão. Ele separa intensidade de duração.

Olhem as calorias por sessão por nível de experiência: o avançado queima 1.265, contra 726 do iniciante, 74\% a mais. A conclusão óbvia seria que quem tem experiência treina mais forte.

Agora, por hora, [\emph{clicar}] os três níveis ficam entre 719 e 722. O avançado não treina mais forte: treina mais tempo, 1,76 hora contra 1,01 do iniciante. Sem essa coluna, eu teria tirado a conclusão errada.

E o contraexemplo: não recriei o IMC, porque a base já traz o IMC calculado a partir de peso e altura. Seria uma coluna repetida, que um dia poderia discordar da original.
\end{fala}
\natela{comece em ``Calorias por sessão'', faça a pergunta, e só então clique em ``Calorias por hora''. As barras mudam na mesma escala, e esse é o momento visual da apresentação.}

% ----------------------------------------------------------------------
\tela{5:40}{7}{O que explica as calorias}
\begin{fala}
Aqui está a correlação de cada coluna com as calorias por sessão. A duração lidera, com 0,91. Experiência e percentual de gordura também aparecem forte, mas pelo mesmo motivo da tela anterior: andam junto com a duração. À direita, numa regressão linear, só a duração explica 82\% da variação; com BPM médio, idade e sexo, chega a 98\%.

Se eu troco o alvo para calorias por hora [\emph{clicar}], a duração some e sobra a intensidade: BPM médio, idade e sexo explicam 95\%, e experiência cai para zero.

Só que 98\% com quatro colunas é alto demais para dado de gente de verdade. Foi isso que me fez desconfiar.
\end{fala}

% ----------------------------------------------------------------------
\tela{6:40}{8}{Explorador}
\begin{fala}
Cada ponto aqui é um praticante. Com BPM médio no eixo X e calorias por hora no Y, os pontos não formam uma nuvem, como seria com pessoas reais: formam quatro retas perfeitas. E à direita, HIIT queima por hora praticamente o mesmo que yoga: a diferença entre os tipos é de 1\%.
\end{fala}
\natela{se sobrar tempo, troque o eixo X para ``Duração da sessão'' e o Y para ``Calorias por sessão'' para mostrar a nuvem normal. Se faltar tempo, diga só a frase das quatro retas e avance.}

% ----------------------------------------------------------------------
\tela{7:20}{9}{A fórmula da base}
\begin{fala}
As quatro retas são uma por combinação de sexo e faixa de idade. Ajustando cada uma, encontrei a regra que gerou a base: calorias igual a duração vezes BPM médio vezes 5,5 para homens ou 5,0 para mulheres, e vezes 0,9 se a pessoa tem mais de 40 anos.

À direita, o real contra o calculado: as 973 linhas caem em cima da diagonal. A maior diferença é de meio kcal, que é só o arredondamento para número inteiro.

Lembram da idade quase plana da tela 3? É mais um sinal: com pessoas reais, a idade não se distribui tão por igual. Isso responde à pergunta de forma exata, mas também mostra que a base é sintética. Peso, tipo de treino e experiência não entram na conta.
\end{fala}
\natela{mexa na calculadora: troque Homem/Mulher e passe a idade de 40 para 41 para mostrar o salto de 10\%.}

\begin{cuidado}
\textbf{Como explicar a descoberta, se perguntarem:} o desvio-padrão da idade (12,18) e do BPM médio (14,35) é praticamente o de uma distribuição uniforme nas mesmas faixas (12,12 e 14,43), o que já sugeria dado gerado. Dividindo calorias por (duração × BPM), sobram só quatro valores: 5,0 e 5,5 até 40 anos, 4,5 e 4,95 acima disso.
\end{cuidado}

% ----------------------------------------------------------------------
\tela{8:40}{10}{Resposta, limitações e próximos passos}
\begin{fala}
Então, a resposta: nesta base, as calorias por sessão dependem de quatro características, que são duração, BPM médio, sexo e idade acima de 40. Peso, tipo de treino, experiência e gordura não entram na conta; quando pareciam importar, era porque andavam junto com a duração.

A principal limitação é que a base é sintética: a resposta descreve o gerador dos dados, não pessoas reais. Isso já está registrado no README como defeito conhecido da fonte. E é uma sessão por pessoa, então não dá para ver a evolução de ninguém.

Os próximos passos são fazer as checagens de valor extremo com IQR e z-score, montar a camada ouro e testar a mesma pergunta numa base real de treinos. Obrigado.
\end{fala}

% ======================================================================
\section{Perguntas prováveis e como responder}

\subsection*{``Se a base é sintética, o projeto perdeu o sentido?''}
Não. O pipeline (ingestão, limpeza, tipos, proveniência) funciona igual para qualquer base, e descobrir que a base é gerada é justamente o papel da exploração e das checagens de qualidade. O erro seria apresentar ``98\% de R²'' como descoberta sobre pessoas.

\subsection*{``Por que a experiência virou categoria e não ficou número?''}
Porque 1, 2 e 3 são níveis, não quantidades: 3 não é o triplo de 1, e a média ``1,8'' não significa nada. Como categoria com ordem, dá para comparar (``mais experiente que'') sem fingir que dá para fazer conta. O número original continua na tabela.

\subsection*{``Por que sexo e tipo de treino como categoria sem ordem?''}
Como texto livre, qualquer erro de digitação vira um grupo novo em silêncio. Como categoria, valor fora da lista vira ausente e aparece na contagem. Sem ordem porque nenhum dos dois tem escala natural: HIIT não é ``maior'' que yoga.

\subsection*{``Por que não removeu os valores extremos?''}
Sinalizar não é remover. Uma sessão fora de 200 a 1.500 kcal/h seria sinalizada, não apagada, e nesta base foram zero. As checagens por IQR e z-score na base de academia são o próximo passo.

\subsection*{``Por que calorias por hora, e não recalcular o IMC?''}
Calorias por hora responde à pergunta: separa intensidade de duração, e as duas colunas de origem existem na base, mas a razão entre elas não. O IMC a base já publica; recriar seria coluna repetida.

\subsection*{``Por que o dicionário de sinônimos do tipo de treino está vazio?''}
Porque a base vem padronizada. Escrever variante que a base não tem seria inventar defeito. O dicionário fica declarado para quando uma extração futura trouxer ``hiit'' ou ``High Intensity''.

\subsection*{``Por que Parquet na prata, e não CSV?''}
O Parquet guarda o tipo de cada coluna: a categoria com ordem continua com ordem quando o arquivo é lido de novo. Num CSV, tudo voltaria a ser texto e número solto.

\subsection*{``Correlação alta quer dizer que a experiência causa gasto calórico?''}
Não, e a tela 6 mostra exatamente isso: a correlação da experiência vem da duração. Correlação não é causa.

% ======================================================================
\section{Antes de apresentar}

\begin{itemize}
  \item Abra o HTML no computador da sala antes da aula e teste as setas. Depois de mexer num controle (lista ou controle deslizante), as setas voltam a trocar de tela sozinhas.
  \item Ensaie duas vezes os dois momentos que dependem de clique: a troca para ``Calorias por hora'' na tela 6 e a calculadora na tela 9.
  \item Deixe o relatório de perfilamento (\texttt{relatorios/}) aberto em outra aba, caso a professora peça para ver a exploração.
\end{itemize}

% ======================================================================
\section{Números de bolso}

\begin{center}
\small
\begin{tabularx}{\linewidth}{@{}X r@{}}
\toprule
\textbf{O quê} & \textbf{Valor}\\
\midrule
Praticantes / colunas na bronze / colunas na prata & 973 / 15 / 18\\
Valores ausentes / linhas repetidas & 0 / 0\\
Homens / mulheres & 511 / 462\\
Força / cardio / yoga / HIIT & 258 / 255 / 239 / 221\\
Iniciante / intermediário / avançado & 376 / 406 / 191\\
Idade (mínimo–máximo, média) & 18–59 anos, 38,7\\
Duração da sessão (mínimo–máximo, média) & 0,5–2 h, 1,26 h\\
Calorias por sessão (mínimo–máximo, média) & 303–1.783 kcal, 905\\
\midrule
\texttt{calorias\_por\_hora}: faixa e média & 540–930 kcal/h, média 720\\
Sessões fora de 200–1.500 kcal/h / durações inválidas & 0 / 0\\
Calorias por sessão: iniciante / intermediário / avançado & 726 / 902 / 1.265\\
Calorias por hora: iniciante / intermediário / avançado & 719 / 722 / 719\\
Duração média: iniciante / avançado & 1,01 h / 1,76 h\\
Intensidade por sexo: feminino / masculino & 684 / 753 kcal/h\\
Diferença máxima entre tipos de treino & 7,5 kcal/h (1\%)\\
\midrule
Correlação com calorias por sessão: duração / experiência / gordura & 0,91 / 0,69 / $-0{,}60$\\
Correlação com calorias por hora: BPM médio / idade & 0,81 / $-0{,}33$\\
R² por sessão: duração / + BPM / + idade / + sexo & 82,5\% / 93,0\% / 95,3\% / 97,9\%\\
R² por hora: BPM / + idade / + sexo / + peso & 65,7\% / 78,9\% / 94,7\% / 94,7\%\\
Linhas que batem com a fórmula (erro máximo) & 973 de 973 (0,5 kcal)\\
\bottomrule
\end{tabularx}
\end{center}

\end{document}
